\documentclass[11pt,a4paper]{article}
\usepackage[T1]{fontenc}
\usepackage[utf8]{inputenc}
\usepackage{lmodern}
\usepackage[left=2.6cm,right=2.6cm,top=2.55cm,bottom=2.1cm,
  headheight=14pt,headsep=20pt,footskip=25pt]{geometry}
\usepackage{amsmath,amssymb,amsthm,mathtools}
\usepackage{booktabs,tabularx,array,enumitem}
\usepackage{fancyhdr,microtype}
\usepackage{xurl}
\usepackage[hidelinks,hypertexnames=false]{hyperref}
\hypersetup{
  pdftitle={A candidate proof of the Gatzweiler-Krattenthaler positivity conjecture},
  pdfauthor={Justin Du},
  pdfsubject={Preprint, version 2; mixed sumsets and the complete parameter range},
  pdfkeywords={Gaussian binomial coefficients, positivity, numerical semigroups,
    Kneser theorem, Hilbert series}}
\setlength{\parindent}{0pt}
\setlength{\parskip}{4pt plus 1pt minus 1pt}
\setlength{\emergencystretch}{1.5em}
\setlist[enumerate]{leftmargin=*,itemsep=3pt,topsep=4pt}
\setlist[itemize]{leftmargin=*,itemsep=3pt,topsep=4pt}
\pagestyle{fancy}
\fancyhf{}
\fancyhead[L]{\small Gatzweiler--Krattenthaler positivity conjecture}
\fancyhead[R]{\small Justin Du -- Preprint v2}
\fancyfoot[C]{\thepage}
\renewcommand{\headrulewidth}{0.3pt}
\renewcommand{\footrulewidth}{0pt}
\numberwithin{equation}{section}
\newtheorem{theorem}{Theorem}[section]
\newtheorem{lemma}[theorem]{Lemma}
\newtheorem{proposition}[theorem]{Proposition}
\newtheorem{corollary}[theorem]{Corollary}
\theoremstyle{remark}
\newtheorem{remark}[theorem]{Remark}
\newcommand{\Z}{\mathbb Z}
\newcommand{\Nzero}{\mathbb Z_{\ge 0}}
\newcommand{\Q}{\mathbb Q}
\newcommand{\C}{\mathbb C}
\newcommand{\sg}[1]{\left\langle #1\right\rangle}
\newcommand{\floor}[1]{\left\lfloor #1\right\rfloor}
\newcommand{\ceil}[1]{\left\lceil #1\right\rceil}
\newcommand{\qbinom}[2]{\genfrac{[}{]}{0pt}{}{#1}{#2}_{q}}
\DeclareMathOperator{\Stab}{Stab}
\DeclareMathOperator{\Hilb}{Hilb}
\newcommand{\Dcond}{\textup{(D)}}
\newcommand{\Hcond}{\textup{(H)}}
\newcommand{\interval}[2]{[#1,#2]_{\Z}}
\newcommand{\code}[1]{\texttt{\small #1}}

\begin{document}
\begin{center}
  {\fontsize{24}{29}\selectfont\bfseries A candidate proof of the\\[4pt]
  Gatzweiler--Krattenthaler\\[4pt]
  positivity conjecture\par}
  \vspace{10pt}
  {\large Justin Du\par}
  \vspace{5pt}
  \href{mailto:fbphdu@gmail.com}{\texttt{fbphdu@gmail.com}}\par
  \vspace{10pt}
  Preprint, version 2 \quad$\vert$\quad 11 November 2026
\end{center}
\vspace{5pt}

\begin{abstract}
Gatzweiler and Krattenthaler conjectured that a polynomial quotient of two
Gaussian binomial coefficients with the same upper index has nonnegative
coefficients. We present a candidate proof for the complete parameter range.
After cancellation, polynomiality gives divisibility-count inequalities for
two separated integer intervals of equal length. We prove that every subset
of the lower interval generates at least as many elements of the upper
interval in its additive semigroup. The criterion of Gatzweiler,
Levic\'an-Santib\'a\~nez and Yoshida then supplies a regular sequence whose
Hilbert series is the required quotient. The central lemma uses a mixed
sumset modulo the largest generator not exceeding the interval length.
Kneser's theorem, together with separate capacities for the lower and upper
generators, handles all stabilizers and all collisions of residues.
Two-generator representations treat subsets near the upper endpoint.
A second sumset estimate handles small lower indices, and a necessary
inequality from polynomiality leaves only two explicit boundary families.
The proof has no finite parameter cutoff; exact computations are supplied
as supplementary checks.
\end{abstract}

{\small\itshape Preprint; not externally peer reviewed.
AI use and verification are described in the appendix.}

\section*{Introduction and proof strategy}

The Gaussian binomial coefficient is a polynomial with nonnegative integer
coefficients, but quotients of such polynomials need not preserve
coefficientwise nonnegativity. Gatzweiler and Krattenthaler~\cite{GK}
proposed that polynomiality is sufficient when the two Gaussian binomial
coefficients have the same upper index. Their paper treats the cases
$l=0,1,2$ and $k-l=0,1,2,3$, and reports exact verification for $n\le 400$.
Gatzweiler, Levic\'an-Santib\'a\~nez and Yoshida~\cite{GLY} subsequently
established a sufficient condition using regular sequences in weighted
polynomial rings. In the normalized range $1\le l<k\le n/2$, their
Corollary~1.7 proves the conjecture when
\[
  2k\floor{\frac{k-1}{2}}+k-l<n,
\]
and hence when $n\ge k^2$.

The present argument uses the semigroup criterion of~\cite{GLY} directly.
It establishes that criterion throughout the normalized parameter range.
Besides the regular-sequence criterion and its Hilbert-series consequence,
the only external theorem used is Kneser's addition theorem, in the form
recorded by DeVos~\cite{DeVos}. All representation estimates and all
parameter reductions needed below are proved in the manuscript.

\textbf{Proof strategy.}
Write $d=k-l$, and cancel the factorials in the Gaussian binomial quotient.
The numerator exponents form an interval $A$ of length $d$ lying above the
denominator interval $B=\interval{l+1}{k}$. Polynomiality says that $A$
contains at least as many multiples of every positive integer as $B$ does.
The task is to show
\[
  |A\cap\sg{I}|\ge |I|
  \qquad(\varnothing\ne I\subseteq B),
\]
where $\sg{I}$ is the additive semigroup generated by $I$.

If $I$ has a generator at most $d$, a mixed sumset yields the desired
inequality. Its two summands record all generators and the generators at
most $d$, respectively. Keeping these two capacities separate is what
allows the periodic and noninjective residue cases to be counted correctly.
If all generators exceed $d$, collisions of direct multiples can instead
be converted into many two-generator representations when $k\ge3l-1$.
A further sumset argument covers $l\le3(\floor{n/k}-1)$. Polynomiality
forces the remaining parameters onto two boundary families, where direct
injections finish the proof.

\textbf{Reading guide.}
Section~\ref{sec:setup} states the results and the interval reduction.
Section~\ref{sec:tools} records the external inputs and an elementary
two-generator bound. Sections~\ref{sec:mixed} and~\ref{sec:tail} prove the
mixed-sumset and upper-endpoint lemmas.
Sections~\ref{sec:three} and~\ref{sec:small} establish the small-index
criterion. Sections~\ref{sec:layers} and~\ref{sec:boundary} identify and
settle the boundary families, and Section~\ref{sec:completion} completes
both the semigroup statement and the original conjecture.

\section{Main results and the interval reduction}\label{sec:setup}

For nonnegative integers $n$ and $0\le j\le n$, put
\begin{equation}\label{eq:qbinom}
  [n]_q!=\prod_{i=1}^{n}\frac{1-q^i}{1-q},
  \qquad [0]_q!=1,
  \qquad
  \qbinom{n}{j}=\frac{[n]_q!}{[j]_q!\,[n-j]_q!}.
\end{equation}
All polynomials and rational functions below have the indeterminate $q$.
The notation $\Nzero[q]$ denotes polynomials with nonnegative integer
coefficients.

\begin{theorem}[The proposed positivity theorem]\label{thm:main}
Let $n,k,l$ be integers with $0\le k,l\le n$. If
\begin{equation}\label{eq:mainquotient}
  \frac{\qbinom{n}{k}}{\qbinom{n}{l}}
\end{equation}
is a polynomial in $q$, then it belongs to $\Nzero[q]$.
\end{theorem}

We will prove a stronger semigroup assertion. For integers $u,v$, write
$\interval{u}{v}=\{x\in\Z:u\le x\le v\}$; this set is empty when $u>v$.
For a finite nonempty set $I$ of positive integers, define
\begin{equation}\label{eq:semigroup}
  \sg{I}=\left\{\sum_{b\in I} t_b b:t_b\in\Nzero\right\}.
\end{equation}
In particular, $0\in\sg{I}$.

\begin{theorem}[Interval semigroup condition]\label{thm:hitting}
Let $0\le l<k\le n/2$, and put
\begin{equation}\label{eq:intervals}
  d=k-l,\qquad N=n-k+1,\qquad U=n-l,\qquad
  A=\interval{N}{U},\qquad B=\interval{l+1}{k}.
\end{equation}
Suppose that, for every positive integer $s$,
\begin{equation}\tag{D}\label{eq:D}
  C_A(s):=|A\cap s\Z|\ \ge\ C_B(s):=|B\cap s\Z|.
\end{equation}
Then, for every nonempty subset $I\subseteq B$,
\begin{equation}\tag{H}\label{eq:H}
  |A\cap\sg{I}|\ge |I|.
\end{equation}
\end{theorem}

The two intervals in~\eqref{eq:intervals} have the same cardinality $d$,
and $N>k$. Factorial cancellation gives
\begin{equation}\label{eq:cancel}
  P_{n,k,l}(q):=
  \frac{\qbinom{n}{k}}{\qbinom{n}{l}}
  =\frac{\prod_{a\in A}(1-q^a)}
         {\prod_{b\in B}(1-q^b)}.
\end{equation}
For $s\ge2$, the exponent of the cyclotomic polynomial $\Phi_s(q)$ in
this rational function is $C_A(s)-C_B(s)$. The factors with $s=1$
cancel because $|A|=|B|$. Consequently
\begin{equation}\label{eq:polynomialcriterion}
  P_{n,k,l}(q)\text{ is a polynomial}
  \quad\Longleftrightarrow\quad \Dcond.
\end{equation}
This also shows that a polynomial quotient automatically lies in
$\Z[q]$: it is a product of cyclotomic polynomials with nonnegative
integer exponents. Thus the phrase \emph{is a polynomial} in
Theorem~\ref{thm:main} does not require an additional integrality
assumption.

For later use, the counts are
\begin{equation}\label{eq:counts}
  C_A(s)=\floor{\frac{U}{s}}-\floor{\frac{N-1}{s}},
  \qquad
  C_B(s)=\floor{\frac{k}{s}}-\floor{\frac{l}{s}}.
\end{equation}
In particular, \Dcond\ implies that every $b\in B$ has a positive
multiple in $A$.

To reduce Theorem~\ref{thm:main} to the normalized range, use
$\qbinom{n}{j}=\qbinom{n}{n-j}$ to move both lower indices into
$\interval{0}{\floor{n/2}}$. Equal indices give the quotient $1$.
The degree $j(n-j)$ is strictly increasing on these integer indices.
If the numerator has the smaller lower index, the rational function has
negative degree and therefore cannot be a nonzero polynomial. All
remaining nontrivial cases have $0\le l<k\le n/2$.

\begin{center}
\small
\begin{tabularx}{\textwidth}{@{}lX@{}}
\toprule
\textbf{Symbol} & \textbf{Meaning in the normalized range}\\
\midrule
$d=k-l$ & Common number of numerator and denominator exponents\\
$A=\interval{N}{U}$ & Numerator exponents, with $N>k$ and $U-N+1=d$\\
$B=\interval{l+1}{k}$ & Denominator exponents\\
$I$ & An arbitrary nonempty subset of $B$\\
$\sg{I}$ & Nonnegative integer combinations of the elements of $I$\\
$t=\floor{n/k}$ & The layer parameter, always at least $2$\\
\bottomrule
\end{tabularx}
\end{center}

\section{Semigroup, sumset, and representation tools}\label{sec:tools}

\subsection{The regular-sequence criterion}

We use the following consequence of Propositions~3.3 and~3.6
of Gatzweiler, Levic\'an-Santib\'a\~nez and Yoshida~\cite{GLY}.

\begin{theorem}[Semigroup criterion]\label{thm:GLY}
Let $A$ and $B$ be finite sets of positive integers with $|A|=|B|$.
If $|A\cap\sg{I}|\ge|I|$ for every nonempty $I\subseteq B$, then
\[
  \frac{\prod_{a\in A}(1-q^a)}{\prod_{b\in B}(1-q^b)}
  \in\Nzero[q].
\]
More precisely, in the weighted polynomial ring
$R=\C[x_b:b\in B]$, with $\deg x_b=b$, there is a homogeneous regular
sequence $(f_a)_{a\in A}$ with $\deg f_a=a$, and the displayed quotient
is $\Hilb(R/(f_a:a\in A);q)$.
\end{theorem}

The coefficient of $q^j$ in this Hilbert series is the dimension of
the degree-$j$ component of the quotient ring. Equal numbers of variables
and regular-sequence elements make that quotient finite dimensional.
This explains both nonnegativity and polynomiality in the conclusion.
Theorem~\ref{thm:hitting}, together with this criterion and
\eqref{eq:polynomialcriterion}, will imply Theorem~\ref{thm:main}.

\subsection{Kneser's theorem and cyclic stabilizers}

For subsets $S,T$ of an abelian group, write
$S+T=\{s+t:s\in S,\ t\in T\}$, and let
$pS=S+\cdots+S$ denote the $p$-fold sumset. The stabilizer of a set $Q$
is $\Stab(Q)=\{h:Q+h=Q\}$.

\begin{theorem}[Kneser]\label{thm:kneser}
Let $S,T$ be nonempty finite subsets of an abelian group, and let
$H=\Stab(S+T)$. Then
\begin{equation}\label{eq:kneser}
  |S+T|\ge |S+H|+|T+H|-|H|.
\end{equation}
\end{theorem}

See~\cite[Theorem~1]{DeVos}. If $S,T,S+T$ occupy respectively $z,w,v$
cosets of $H$, then $S+T$ is a union of full $H$-cosets and
\begin{equation}\label{eq:knesercosets}
  v\ge z+w-1.
\end{equation}
For any positive integer $p$ and any nonempty finite set $S$, we will also use
\begin{equation}\label{eq:kneseriterated}
  |pS|\ge p|S+H|-(p-1)|H|,
  \qquad H=\Stab(pS).
\end{equation}
To justify this form, pass to the quotient by $H$.
The image of $pS$ has trivial stabilizer. The stabilizer of every
intermediate sumset $jS$ in the quotient is contained in that of $pS$,
because adding $(p-j)S$ preserves every period of $jS$.
Thus each intermediate sumset is aperiodic. Applying
\eqref{eq:kneser} successively in the quotient gives
$|p\overline S|\ge p|\overline S|-(p-1)$, which lifts to
\eqref{eq:kneseriterated}.

Every subgroup of $\Z/m\Z$ has the form
\begin{equation}\label{eq:cyclicsubgroup}
  H=\{0,g,2g,\ldots,m-g\},\qquad g\mid m.
\end{equation}
Its order is $h=m/g$, and its cosets are indexed by residues modulo $g$.
A set stabilized by $H$ therefore contains every residue modulo $m$
lying in any of the modulo-$g$ classes it occupies.

\subsection{An elementary two-generator bound}

\begin{lemma}\label{lem:frobenius}
Let $a,b$ be distinct positive integers. Put
$g=\gcd(a,b)$, $L=ab/g$, and $F=L-a-b$.
Every nonnegative multiple of $g$ strictly greater than $F$ belongs
to $\sg{a,b}$.
\end{lemma}

\begin{proof}
Write $a=ug$ and $b=vg$, with $\gcd(u,v)=1$.
For an integer $w>uv-u-v$, choose
$j\in\{0,1,\ldots,u-1\}$ such that $jv\equiv w\pmod u$.
Then
\[
  w-jv>uv-u-v-(u-1)v=-u.
\]
Since $w-jv$ is a multiple of $u$, it is nonnegative.
Hence $w=(w-jv)+jv\in\sg{u,v}$.
Scaling by $g$ proves the assertion.
\end{proof}



\section{The mixed-sumset lemma}\label{sec:mixed}

Throughout this section, $A$ and $B$ are as in
\eqref{eq:intervals}, with $N>k$ and \Dcond.
The following lemma treats every subset containing a generator no
larger than the common interval length.

\begin{lemma}[Mixed sumsets]\label{lem:mixed}
If $\varnothing\ne I\subseteq B$ and $I\cap\interval{1}{d}\ne\varnothing$,
then $|A\cap\sg{I}|\ge|I|$.
\end{lemma}

\begin{proof}
Define
\begin{equation}\label{eq:mixedlow}
  I_{\mathrm{low}}=I\cap\interval{1}{d},\qquad
  m=\max I_{\mathrm{low}},\qquad M=\max I.
\end{equation}
Then $l<m\le d$, and
\begin{equation}\label{eq:mixedranges}
  I_{\mathrm{low}}\subseteq\interval{l+1}{m},
  \qquad
  I\setminus I_{\mathrm{low}}\subseteq\interval{d+1}{k}.
\end{equation}
The second interval has length $l$. There are no elements of $I$ in
$\interval{m+1}{d}$; this uses the maximal choice of $m$.

\textbf{Residues and their lifts.}
In $\Z/m\Z$, put
\begin{equation}\label{eq:mixedsumset}
  S=I\bmod m,\qquad T=I_{\mathrm{low}}\bmod m,\qquad Q=S+T.
\end{equation}
Both $S$ and $T$ contain $0$, since $m\in I_{\mathrm{low}}$.
Suppose $y\in A$ and $y\bmod m\in Q$.
Represent its residue as a sum of one residue from $S$ and one from $T$.
For a nonzero residue choose a corresponding generator, and for a zero
residue choose the empty sum. The resulting integer $\zeta$ belongs to
$\sg{I}$ and satisfies
\begin{equation}\label{eq:mixedlift}
  0\le \zeta\le M+m,\qquad \zeta\equiv y\pmod m.
\end{equation}
If $\zeta>y$, the congruence would give
$\zeta\ge y+m\ge N+m>M+m$, a contradiction.
Therefore $y=\zeta+jm$ for some $j\ge0$, and $y\in\sg{I}$.
It suffices to count the elements of $A$ whose residues lie in $Q$.

\textbf{Capacities of occupied cosets.}
Let $H=\Stab(Q)$, write it as in~\eqref{eq:cyclicsubgroup}, and put
$h=m/g$. Let $S,T,Q$ occupy respectively $z,w,v$ cosets of $H$.
These are the occupied classes modulo $g$, and Kneser's theorem gives
\begin{equation}\label{eq:mixedkneser}
  v\ge z+w-1.
\end{equation}
Set
\begin{equation}\label{eq:alphabeta}
  \alpha=\floor{\frac{l}{g}},\qquad
  \beta=\ceil{\frac{l}{g}}.
\end{equation}
Since $l<m=gh$, we have $0\le\alpha\le\beta\le h$.

Count the elements of $I$ by their classes modulo $g$.
The zero class contains at most $C_B(g)$ elements of $I$ in total.
This count includes both parts of~\eqref{eq:mixedranges}.
In each nonzero class, $\interval{l+1}{m}$ contains at most
$h-\alpha$ elements, since
\[
  \ceil{\frac{m-l}{g}}=h-\floor{\frac{l}{g}}.
\]
There are at most $w-1$ nonzero classes occupied by the low generators.
The high interval $\interval{d+1}{k}$ has length $l$ and contains
at most $\beta$ elements in any class; at most $z-1$ nonzero classes
are relevant for its elements. Hence
\begin{equation}\label{eq:mixedupper}
  |I|\le C_B(g)+(w-1)(h-\alpha)+(z-1)\beta.
\end{equation}

On the other hand, $0\in Q$ and $Q$ is $H$-periodic, so the entire
zero coset $H$ is contained in $Q$. In $A$ this coset contributes
$C_A(g)\ge C_B(g)$ elements by \Dcond.
Each other occupied coset contributes at least
$\floor{d/g}\ge h$ elements, because $d\ge m=gh$.
All of these elements are representable by~\eqref{eq:mixedlift}.
Consequently
\begin{equation}\label{eq:mixedlower}
  |A\cap\sg{I}|
  \ge C_B(g)+(v-1)h
  \ge C_B(g)+(z+w-2)h.
\end{equation}
The last lower bound minus the upper bound in~\eqref{eq:mixedupper}
equals
\begin{equation}\label{eq:mixeddifference}
  (z-1)(h-\beta)+(w-1)\alpha\ge0.
\end{equation}
This proves the lemma.
\end{proof}

\begin{remark}\label{rem:mixedlocal}
The proof does not require the map $I\longrightarrow\Z/m\Z$ to be
injective. It counts original generators through their interval capacities,
while Kneser's theorem counts occupied cosets. It also allows a nontrivial
stabilizer. In fact, after the stabilizer modulus $g$ has been determined,
the proof uses \Dcond\ only for that one modulus.
\end{remark}

\section{Subsets near the upper endpoint}\label{sec:tail}

We next handle subsets having no generator at most $d$, provided that
$k\ge3l-1$. This lemma needs only the existence of a direct multiple
for each generator.

\begin{lemma}[Upper-endpoint subsets]\label{lem:tail}
Let $l\ge1$, $k\ge3l-1$, and $d=k-l$.
Let $A=\interval{N}{U}$ have length $d$, with $N>k$.
Suppose $\varnothing\ne I\subseteq\interval{d+1}{k}$ and every $b\in I$
has a multiple in $A$. Then $|A\cap\sg{I}|\ge|I|$.
\end{lemma}

\begin{proof}
Let $D_I$ be the union, over $b\in I$, of the multiples of $b$ in $A$.
If $|D_I|\ge|I|$, the conclusion is immediate.
Otherwise, since every generator has at least one such multiple, two
distinct generators $a<b$ have a common multiple $x\in A$.
Put
\begin{equation}\label{eq:tailparameters}
  g=\gcd(a,b),\quad a=ug,\quad b=vg,\quad
  J=\floor{\frac{k}{g}},\quad
  L=\frac{ab}{g},\quad F=L-a-b.
\end{equation}
Since $x$ is a positive common multiple, $U\ge x\ge L$.

The interval $\interval{d+1}{k}$ has length $l$. Thus
$g\le b-a\le l-1$, and
\[
  d=k-l\ge2l-1\ge l+g.
\]
It follows that
\begin{equation}\label{eq:tailwidth}
  \floor{\frac{d}{g}}\ge\ceil{\frac{l}{g}}.
\end{equation}
For example, to see the floor-ceiling step, use
$d/g\ge l/g+1$ and
$\floor{l/g}+1\ge\ceil{l/g}$.
The collision case is impossible for $l=1$, so all the expressions
above are consistent with $g\ge1$.

Since $k\ge3l-1\ge2l$, we have $d=k-l\ge k/2$.
Thus $u,v>k/(2g)\ge J/2$.
The distinct integers $u<v$ therefore satisfy
\begin{equation}\label{eq:tailsum}
  u+v\ge J+2.
\end{equation}
Indeed, if $J=2j$, their sum is at least $2j+3$; if $J=2j+1$,
it is at least $2j+3$.

Fix a residue class modulo $g$ occupied by $I$.
Let $c$ be the smallest generator in this class, and let $r_c$ be its
number of generators. We have both
\begin{equation}\label{eq:tailrc}
  r_c\le\ceil{\frac{l}{g}},
  \qquad
  r_c\le\floor{\frac{k-c}{g}}+1.
\end{equation}
We will produce $r_c$ points $y\in A$ satisfying
$y\equiv c\pmod g$ and $y>F+c$.
Since $y>k\ge c$, the difference $y-c$ is nonnegative.
Lemma~\ref{lem:frobenius} then gives
\[
  y\in c+\sg{a,b}\subseteq\sg{I}.
\]

Define
\begin{equation}\label{eq:tailcapacities}
  \begin{split}
  C_c&=\floor{\frac{U-c}{g}}-\floor{\frac{N-1-c}{g}},\\
  Q_c&=\floor{\frac{U-c}{g}}-\frac{F}{g}.
  \end{split}
\end{equation}
Because $F/g$ is an integer, the exact number of selected points is
\begin{equation}\label{eq:tailmin}
  \max\{0,\min(C_c,Q_c)\}.
\end{equation}
For completeness, put
$j_{\max}=\floor{(U-c)/g}$ and
$j_{\min}=\floor{(N-1-c)/g}+1$.
The points in $A$ in this residue class are $c+jg$ for
$j_{\min}\le j\le j_{\max}$.
The threshold $y>F+c$ imposes $j\ge F/g+1$.
Thus their number is
$\max\{0,j_{\max}-\max(j_{\min},F/g+1)+1\}$,
which is exactly~\eqref{eq:tailmin}.

Every residue class has at least $\floor{d/g}$ points in an interval
of length $d$. Equations~\eqref{eq:tailwidth} and~\eqref{eq:tailrc}
give $C_c\ge r_c$.
Using $U\ge L$ and~\eqref{eq:tailsum}, we also obtain
\begin{equation}\label{eq:tailQbound}
  \begin{split}
  Q_c
  &\ge u+v-\ceil{\frac{c}{g}}\\
  &\ge J+2-\ceil{\frac{c}{g}}\\
  &\ge \floor{\frac{k-c}{g}}+1
   \ge r_c.
  \end{split}
\end{equation}
The penultimate inequality follows from
\[
  \ceil{\frac{c}{g}}+\floor{\frac{k-c}{g}}\le J+1.
\]
For an explicit calculation, write $k=Jg+r$ and $c=pg+s$, with
$0\le r,s<g$. Then
\[
  \ceil{\frac{c}{g}}+\floor{\frac{k-c}{g}}
  =J+\mathbf{1}_{s>0}+\floor{\frac{r-s}{g}}\le J+1.
\]
Here $\mathbf{1}_{s>0}$ is $1$ when $s>0$ and $0$ otherwise.
Hence~\eqref{eq:tailmin} is at least $r_c$.
Summing over the occupied residue classes proves the lemma, since those
classes are disjoint.
\end{proof}

\begin{corollary}\label{cor:tailrange}
Under the hypotheses of Theorem~\ref{thm:hitting}, if $k\ge3l-1$,
then \Hcond\ holds for every nonempty $I\subseteq B$.
\end{corollary}

\begin{proof}
If $I$ has an element at most $d$, apply Lemma~\ref{lem:mixed}.
Otherwise $I\subseteq\interval{d+1}{k}$.
If $l=0$, this latter case is empty; if $l\ge1$, \Dcond\ supplies the
direct multiples required by Lemma~\ref{lem:tail}.
\end{proof}

\begin{remark}
The modulus $g$ in Lemma~\ref{lem:tail} is the greatest common divisor
of one collision pair. No divisibility by $g$ is required of the other
generators. Each occupied class is treated through its own translate
$c+\sg{a,b}$, and~\eqref{eq:tailmin} keeps the left endpoint $N$
in the count.
\end{remark}

\section{At most three generators}\label{sec:three}

This section proves a small-cardinality statement for a more general
pair of positive integer intervals. It will supply the initial cases
of the iterated-sumset argument.

\begin{lemma}\label{lem:three}
Let $A,B$ be nonempty finite intervals of positive integers, with
$\min A>\max B$ and $|A|\ge|B|$.
Assume $|A\cap s\Z|\ge|B\cap s\Z|$ for every positive integer $s$.
Then $|A\cap\sg{I}|\ge|I|$ for every $I\subseteq B$ with
$1\le|I|\le3$.
\end{lemma}

\begin{proof}
\textbf{One or two generators.}
For $I=\{b\}$, the multiple-count inequality with $s=b$ gives a
representable point in $A$.
For $I=\{a,b\}$, consider the union of their direct multiples in $A$.
If it has at least two elements, there is nothing to prove.
Otherwise both generators have the same unique direct multiple $x\in A$.
Let $g=\gcd(a,b)$ and $L=ab/g$.
The interval $B$ has at least two multiples of $g$, so $A$ does as well.
Since the multiples of $g$ in an interval are consecutive along the
$g$-lattice, at least one of $x-g$ and $x+g$ lies in $A$.
That point is positive and at least
\[
  x-g\ge L-g>L-a-b.
\]
It is representable by Lemma~\ref{lem:frobenius}, and is distinct
from $x$.

\textbf{Three generators.}
Suppose $|I|=3$.
If the union of all direct multiples has at least three points,
we are done. Otherwise there is a collision pair $a<b$ in $I$ with
a common multiple $x\in A$; write $c$ for the third generator, and put
\[
  g=\gcd(a,b),\qquad L=ab/g,\qquad F=L-a-b.
\]
If $A$ contains at least three multiples of $g$, choose three
consecutive such multiples containing $x$.
Their smallest is at least $x-2g\ge L-2g>F$, since $a+b>2g$.
All three are representable by Lemma~\ref{lem:frobenius}.

Since $a$ and $b$ are distinct multiples of $g$ in $B$, the
multiple-count hypothesis gives
\[
  2\le |B\cap g\Z|\le |A\cap g\Z|.
\]
The preceding case covers $|A\cap g\Z|\ge3$. Thus the only remaining
case has exactly two multiples of $g$ in each of $A$ and $B$.
Consequently, for some integer $v\ge1$,
\begin{equation}\label{eq:threeconsecutive}
  a=vg,\qquad b=(v+1)g,\qquad c<(v+2)g.
\end{equation}
The first two expressions follow because two multiples in an interval
with no third must be consecutive. The last inequality holds because
$B$ does not contain the next multiple $(v+2)g$.
The only multiples of $g$ in $B$ are $a$ and $b$.
Since $c\in B\setminus\{a,b\}$, we have $g\nmid c$, and consequently
$g>1$. These conclusions follow from the multiple count, with no
assumption on the overall greatest common divisor of $I$.
We also have
\begin{equation}\label{eq:threelength}
  |A|\ge|B|\ge g+1,\qquad
  \min A>x-2g\ge L-2g.
\end{equation}
The last strict inequality holds because $x-2g$ in $A$ would give
three multiples of $g$.
The two multiples of $g$ in $A$ are themselves representable:
they are at least $x-g\ge L-g>F$.
We seek one further point in the nonzero residue class of $c$.

\textbf{The remaining cases.}
If $v=1$, then $a=g$.
By~\eqref{eq:threelength}, $A$ contains a point congruent to $c$
modulo $g$; since that point exceeds $c$, it belongs to
$c+\sg{g}$.

If $v\ge3$, then~\eqref{eq:threeconsecutive} gives
\[
  F+c<L-(v-1)g\le L-2g<\min A.
\]
Every point of $A$ in the residue class of $c$ therefore belongs to
$c+\sg{a,b}$. At least one exists by~\eqref{eq:threelength}.

Suppose now $v=2$. Then $a=2g$, $b=3g$, $L=6g$, and $F=g$.
If $c\le3g$, we have $F+c\le4g<\min A$, and the same argument
applies.
If $c>3g$, write $c=3g+s$ with $1\le s<g$, and $x=6jg$
with $j\ge1$. Both integers
\begin{equation}\label{eq:threewitnesses}
  \begin{split}
  y&=x-g+s=c+2g(3j-2),\\
  z&=x+s=c+3g(2j-1)
  \end{split}
\end{equation}
belong to $\sg{I}$ and have the nonzero residue of $c$.
They satisfy $y<x<z$ and $z-y=g$.
An integer interval containing $x$ and avoiding both $y$ and $z$
would be contained in $\interval{y+1}{z-1}$ and have length at most
$g-1$. This contradicts $|A|\ge g+1$.
Thus at least one of $y,z$ belongs to $A$ and supplies the third
representable point. All cases are complete.
\end{proof}

\section{The small-index criterion}\label{sec:small}

\begin{lemma}\label{lem:small}
Under the hypotheses of Theorem~\ref{thm:hitting}, let
$t=\floor{n/k}\ge2$. If $l\le3(t-1)$, then \Hcond\ holds for every
nonempty $I\subseteq B$.
\end{lemma}

\begin{proof}
If $|I|\le3$, apply Lemma~\ref{lem:three}.
Henceforth put $r=|I|\ge4$ and $m=\max I$.
If $m\le d$, move each generator $b\in I$ forward by a nonnegative
multiple of $m$ into the interval $\interval{N}{N+m-1}$.
The resulting points are representable and distinct: different
elements of $\interval{1}{m}$ have different residues modulo $m$.
Since $m\le d$, the whole target interval lies in $A$.

We may therefore suppose $m>d$. Put
\[
  e=m-d\le l,\qquad
  S=I\bmod m,\qquad W=A\bmod m.
\]
Then $|S|=r$, $|W|=d$, and $0\in S$.
The equalities of cardinalities use $I\subseteq\interval{1}{m}$
and $d<m$, respectively.

\textbf{Lifting the iterated sumset.}
Since $n\ge tk$,
\begin{equation}\label{eq:smallN}
  N=n-k+1>(t-1)k\ge(t-1)m.
\end{equation}
If $y\in A$ and $y\bmod m\in tS$, choose at most $t$ generators
representing its summands, using the empty sum for a zero residue.
Every chosen nonzero generator is less than $m$, so their sum $\zeta$
satisfies $0\le \zeta<tm$.
If $\zeta>y$, congruence modulo $m$ and~\eqref{eq:smallN} would give
$\zeta\ge y+m>tm$, a contradiction.
Thus $y=\zeta+jm\in\sg{I}$ for some $j\ge0$.
It remains to show $|W\cap tS|\ge r$.

\textbf{Counting with the stabilizer.}
Write $H=\Stab(tS)$ in the form~\eqref{eq:cyclicsubgroup}, and put
$h=m/g$. Let $S$ and $tS$ occupy $z$ and $w$ cosets of $H$.
By~\eqref{eq:kneseriterated},
\begin{equation}\label{eq:smallkneser}
  w\ge tz-(t-1)\ge2z-1.
\end{equation}
The second inequality uses $t\ge2$ and $z\ge1$.
Define
\[
  v=h-\floor{\frac{l}{g}}.
\]
Because $l<m=gh$, we have $v\ge1$.
The interval $\interval{l+1}{m}$ contains at most $v$ integers
in each class modulo $g$, so
\begin{equation}\label{eq:smallcapacity}
  r\le zv.
\end{equation}
Every coset contains at least $v-1$ elements of $W$.
Indeed, $d\ge m-l$ implies
\[
  \floor{\frac{d}{g}}
  \ge h-\ceil{\frac{l}{g}}
  \ge v-1.
\]
For the zero coset, \Dcond\ gives the stronger bound
\begin{equation}\label{eq:smallzero}
  |W\cap H|=C_A(g)\ge C_B(g)
  =\floor{\frac{k}{g}}-\floor{\frac{l}{g}}\ge v,
\end{equation}
since $k\ge m$.
The equality $|W\cap H|=C_A(g)$ uses the injectivity of reduction
of $A$ modulo $m$.

If $v\ge2$, then $0\in tS$ implies $H\subseteq tS$.
Using~\eqref{eq:smallkneser} and~\eqref{eq:smallzero}, we get
\begin{equation}\label{eq:smallvlarge}
  |W\cap tS|
  \ge v+(2z-2)(v-1)
  =zv+(z-1)(v-2)\ge r.
\end{equation}

Suppose $v=1$. Each occupied coset of $S$ now contains exactly one
element of $S$, so $z=r$.
If $d\ge g$, every coset has at least one element of $W$; hence
$|W\cap tS|\ge w\ge r$.
If $d<g$, the interval $A$ occupies $d$ distinct classes modulo $g$.
There are $g$ such classes in total, and $tS$ occupies $w$ of them.
Therefore
\begin{equation}\label{eq:smallvone}
  |W\cap tS|\ge d+w-g
  \ge tr-(t-1)-(g-d).
\end{equation}
Finally,
\[
  (t-1)(r-1)\ge3(t-1)\ge l\ge m-d\ge g-d,
\]
where we used $r\ge4$, the hypothesis on $l$, and $m\le k$.
Thus the last expression in~\eqref{eq:smallvone} is at least $r$.
This proves the lemma.
\end{proof}

\section{Layer restrictions forced by polynomiality}\label{sec:layers}

The next proposition uses only the multiple-count condition.
It is the parameter reduction that makes the preceding two ranges
sufficient, up to two explicit boundary families.

\begin{proposition}\label{prop:layers}
Under the hypotheses of Theorem~\ref{thm:hitting}, suppose $d\ge2$
and put $t=\floor{n/k}\ge2$. Then
\begin{equation}\label{eq:layerk}
  k\ge(t+1)(l-t+1),
\end{equation}
and
\begin{equation}\label{eq:layern}
  tk+l\le n\le
  \min\{(t+1)k-1,\ (t+1)k-tl+t^2-1\}.
\end{equation}
\end{proposition}

\begin{proof}
The definition of $t$ gives
\[
  (t-1)k<N\le tk,\qquad U<(t+1)k.
\]
Since $k\in B$, condition \Dcond\ forces a multiple of $k$ in $A$.
The only possible one is $tk$, so $tk\in A$.
As $A$ has length $d$,
\begin{equation}\label{eq:layerNlower}
  N\ge tk-d+1.
\end{equation}

Now take
\[
  b=\floor{\frac{N-1}{t}}.
\]
Since $N\le tk$, we have $b\le k-1$.
We claim $b\ge l+1$.
Otherwise $b\le l$ implies $N\le t(l+1)$.
Both $t(l+1)$ and $tk$ would then lie in $A$: they are between $N$
and $tk\le U$.
Their difference is $t(d-1)>d-1$, which is impossible in an interval
of length $d$. Thus $b\in B$.

By construction $tb<N$.
Condition \Dcond\ supplies a multiple of $b$ in $A$, and its multiplier
must be at least $t+1$. Hence
\[
  (t+1)b\le U=N+d-1.
\]
Write $N=tb+v$ with $1\le v\le t$.
The last inequality yields $b\le d+v-1\le d+t-1$, and therefore
\begin{equation}\label{eq:layerNupper}
  N\le t(d+t).
\end{equation}
Combining~\eqref{eq:layerNlower} and~\eqref{eq:layerNupper}, and using
$d=k-l$, gives
\[
  tk-(k-l)+1\le t(k-l+t),
  \qquad
  k\ge(t+1)l-t^2+1=(t+1)(l-t+1).
\]
This proves~\eqref{eq:layerk}.
Adding $k-1$ to~\eqref{eq:layerNlower} and~\eqref{eq:layerNupper}
gives the corresponding lower and upper bounds for $n$.
The further bound $n\le(t+1)k-1$ is part of the definition of $t$.
Together these are~\eqref{eq:layern}.
\end{proof}

\section{The two boundary families}\label{sec:boundary}

The relevant boundary parameters have $t=2$ and $k$ equal to
$3l-3$ or $3l-2$. We first determine $n$, and then give explicit
injections for all subsets having no generator at most $d$.

\begin{lemma}\label{lem:boundaryparameters}
Assume the hypotheses of Theorem~\ref{thm:hitting}, with $l\ge3$ and
$t=2$.
\begin{enumerate}[label=\textup{(\roman*)}]
\item If $k=3l-3$, then $n=7l-6$.
\item If $k=3l-2$, then $n=7l-3$.
\end{enumerate}
\end{lemma}

\begin{proof}
For $t=2$, equation~\eqref{eq:layern} gives
\[
  2k+l\le n\le3k-2l+3.
\]
For $k=3l-3$, the two bounds are both $7l-6$.
For $k=3l-2$, they give $n\in\{7l-4,7l-3\}$.
The first choice violates \Dcond.
Indeed, $b=2l-1$ belongs to $B$, whereas for $n=7l-4$ one has
\[
  A=\interval{4l-1}{6l-4},\qquad
  2b=4l-2<\min A,\qquad 3b=6l-3>\max A.
\]
Since two consecutive multiples of $b$ straddle the whole interval,
$A$ contains no multiple of $b$.
Thus only $n=7l-3$ is possible.
\end{proof}

\begin{lemma}[Boundary injections]\label{lem:boundaryinjections}
Let $l\ge3$.
For either of the two parameter families
\[
  (k,n)=(3l-3,7l-6)
  \quad\text{or}\quad
  (k,n)=(3l-2,7l-3),
\]
every nonempty $I\subseteq\interval{d+1}{k}$ satisfies
$|A\cap\sg{I}|\ge|I|$.
This assertion does not require \Dcond.
\end{lemma}

\begin{proof}
\textbf{The family $k=3l-3$.}
Set $g=l-1\ge2$.
The parameters become
\[
  k=3g,\qquad d=2g-1,\qquad
  A=\interval{4g+2}{6g},\qquad
  I\subseteq\interval{2g}{3g}.
\]
If $3g\in I$, the map
\begin{equation}\label{eq:boundaryonea}
  c\longmapsto c+3g\qquad(c\in I)
\end{equation}
is injective and has representable values in
$\interval{5g}{6g}\subseteq A$.
The inclusion uses $g\ge2$.
If $3g\notin I$, use
\begin{equation}\label{eq:boundaryoneb}
  c\longmapsto
  \begin{cases}
    6g,&c=2g,\\
    2c,&c>2g.
  \end{cases}
\end{equation}
The exceptional value is $3(2g)$.
Every other value is an even integer between $4g+2$ and $6g-2$.
They are mutually distinct and avoid $6g$.
Thus~\eqref{eq:boundaryoneb} is also an injection into $A\cap\sg{I}$.

\textbf{The family $k=3l-2$.}
Here
\[
  d=2l-2,\qquad
  A=\interval{4l}{6l-3},\qquad
  I\subseteq\interval{2l-1}{3l-2}.
\]
The map
\begin{equation}\label{eq:boundarytwo}
  c\longmapsto
  \begin{cases}
    3c=6l-3,&c=2l-1,\\
    2c,&c\ge2l
  \end{cases}
\end{equation}
has representable values in $A$.
Its exceptional value is odd, whereas the other values are distinct
even integers in $\interval{4l}{6l-4}$.
This proves injectivity and completes both families.
\end{proof}

\section{Completion of the proof}\label{sec:completion}

\begin{proof}[Proof of Theorem~\ref{thm:hitting}]
If $d=1$, then $B=\{k\}$ and $A=\{n-k+1\}$.
Condition \Dcond\ with $s=k$ gives $k\mid n-k+1$, so
$A\subseteq\sg{k}$ and \Hcond\ follows.
The corresponding quotient is the finite geometric sum
\[
  \frac{1-q^{n-k+1}}{1-q^k}
  =\sum_{j=0}^{(n-k+1)/k-1}q^{jk}.
\]

Suppose $d\ge2$, and put $t=\floor{n/k}\ge2$.
If $l\le3(t-1)$, Lemma~\ref{lem:small} proves \Hcond\ for every
nonempty subset of $B$.
If $k\ge3l-1$, Corollary~\ref{cor:tailrange} gives the same conclusion.
It remains to suppose that both inequalities fail:
\begin{equation}\label{eq:remainingrange}
  l\ge3t-2,\qquad k<3l-1.
\end{equation}
Proposition~\ref{prop:layers} gives
\begin{equation}\label{eq:completiondifference}
  \begin{split}
  k-(3l-1)
  &\ge (t-2)l-t^2+2\\
  &\ge (t-2)(3t-2)-t^2+2\\
  &=2(t-1)(t-3).
  \end{split}
\end{equation}
The second inequality uses $t\ge2$.
If $t\ge3$, the last expression is nonnegative, contradicting
$k<3l-1$. Hence $t=2$.
Equation~\eqref{eq:layerk} now says $k\ge3l-3$, so the only
integer choices are
\begin{equation}\label{eq:onlyboundaries}
  k=3l-3\quad\text{or}\quad k=3l-2.
\end{equation}
Moreover $l\ge4$ by~\eqref{eq:remainingrange}.
Lemma~\ref{lem:boundaryparameters} supplies the corresponding values
of $n$.

Fix any nonempty $I\subseteq B$ in either remaining family.
If it contains a generator at most $d$, apply Lemma~\ref{lem:mixed}.
Otherwise apply Lemma~\ref{lem:boundaryinjections}.
Thus every such $I$ satisfies \Hcond.
The cases above cover all normalized parameters, proving
Theorem~\ref{thm:hitting}.
\end{proof}

\begin{center}
\small
\begin{tabularx}{\textwidth}{@{}p{0.43\textwidth}X@{}}
\toprule
\textbf{Parameter range} & \textbf{Argument giving \Hcond\ for all subsets}\\
\midrule
$d=1$ & The unique denominator generator divides the numerator exponent\\
$d\ge2,\ l\le3(t-1)$ & Lemma~\ref{lem:small}\\
$d\ge2,\ k\ge3l-1$ & Mixed sumsets and the upper-endpoint lemma\\
Neither preceding inequality, $d\ge2$ &
$t=2$ and the two boundary families;
mixed sumsets and explicit injections\\
\bottomrule
\end{tabularx}
\end{center}

\begin{proof}[Proof of Theorem~\ref{thm:main}]
Reduce to $0\le l<k\le n/2$ as in Section~\ref{sec:setup}.
Polynomiality gives \Dcond\ by~\eqref{eq:polynomialcriterion}.
Theorem~\ref{thm:hitting} gives \Hcond\ for every nonempty
$I\subseteq B$.
Applying Theorem~\ref{thm:GLY} to~\eqref{eq:cancel} yields
$P_{n,k,l}(q)\in\Nzero[q]$.
The equal-index cases have quotient $1$, and the other ordering
cannot give a polynomial. This proves the full statement.
\end{proof}

\begin{corollary}[A Hilbert-series realization]\label{cor:hilbert}
Every polynomial quotient in the normalized range of
Theorem~\ref{thm:main} is the Hilbert series of a finite-dimensional
graded complete intersection
\[
  \C[x_{l+1},\ldots,x_k]/(f_N,\ldots,f_U),
  \qquad \deg x_b=b,\quad \deg f_a=a.
\]
\end{corollary}

\begin{proof}
Theorem~\ref{thm:hitting} verifies the hypothesis of the
regular-sequence assertion in Theorem~\ref{thm:GLY}.
\end{proof}

\section*{Appendix. AI use and supplementary verification}
\addcontentsline{toc}{section}{Appendix. AI use and supplementary verification}

\textbf{Changes in version 2.}
The two-multiple branch in Lemma~\ref{lem:three} now derives
$g\nmid c$ and $g>1$ directly from the count in $B$; the unnecessary
overall-coprimality reduction has been removed.
The right-endpoint floor-ceiling bound is expanded explicitly, and the
lifted integers in the sumset arguments use a different symbol from
coset counts. Reproduction commands preserve literal double hyphens.
The release-check metadata and checksum manifest now refer only to files
actually included in the package.

\textbf{Author and AI use.}
Justin Du supplied the candidate proof and verification materials and
directed the preparation of this preprint.
The manuscript was prepared in an AI-assisted workflow.
OpenAI Codex was used to review the supplied argument, expand its
mathematical exposition, check references, rerun the verification programs,
and prepare the \LaTeX\ source and PDF.
The author is responsible for the preprint and is the contact for
questions and corrections.

\textbf{Review status.}
The argument has undergone internal AI-assisted mathematical review
and exact-arithmetic checks. External human peer review and
proof-assistant formalization are not supplied with this version.
The word \emph{candidate} records that review status.
All parameter branches are addressed in the text, relative to the two
external inputs explicitly identified in Section~\ref{sec:tools}.
This version makes no claim of bibliographic priority.

\textbf{Role of the computations.}
The proof of Theorems~\ref{thm:hitting} and~\ref{thm:main} is the
symbolic argument in Sections~\ref{sec:setup}--\ref{sec:completion}.
No computational cutoff or random sample is used as a premise.
The supplied programs test the constructions, capacity inequalities,
and exact coefficient identities on finite sets of inputs.
Algorithmic cross-checks are distinct from external mathematical review.

\textbf{Main audit.}
The program \code{verify\_mixed\_completion.py} scans every triple
$(k,l,n)$ with
\[
  2\le k\le16,\qquad 1\le l<k,\qquad 2k\le n\le7k.
\]
It tests \Dcond\ and, for each polynomial instance, checks every
nonempty subset of $B$.
The randomized stage accepts $6{,}000$ polynomial parameter draws
with $15\le k\le600$ and tests selected subsets, including subsets
with few low generators and subsets confined to the upper endpoint.
These draws need not be distinct.
The two boundary injections are also checked for $3\le l\le300$
without imposing polynomiality.
The coefficient checks multiply each selected full quotient back by
the denominator and compare with the complete numerator polynomial,
without truncating the product identity.

\begin{center}
\small
\begin{tabularx}{\textwidth}{@{}Xr@{}}
\toprule
\textbf{Main audit item} & \textbf{Exact count}\\
\midrule
Parameter triples in the exhaustive scan & $6{,}920$\\
Polynomial triples in that scan & $971$\\
Random accepted polynomial draws & $6{,}000$\\
Subset checks, including selected random subsets & $3{,}745{,}276$\\
Mixed-sumset checks & $3{,}726{,}163$\\
Mixed checks with a nontrivial stabilizer & $2{,}399{,}416$\\
Mixed checks with collisions under reduction modulo $m$ & $517{,}086$\\
Comparisons of two semigroup algorithms & $1{,}494$\\
Boundary-injection checks without polynomiality & $4{,}768$\\
Complete polynomial identities & $58$\\
Quotient coefficients in those identities & $82{,}340$\\
\bottomrule
\end{tabularx}
\end{center}

The main semigroup oracle uses a bitset and successive doubling of the
available number of each generator.
A second implementation computes membership by the recurrence
$u_0=1$ and
$u_y=\max_{b\in I,\ b\le y}u_{y-b}$, taking the maximum of an empty
set to be $0$. The $1{,}494$ comparisons between these implementations
all agree.
For a constructed certificate, the program checks that its points
actually belong to the independently computed semigroup.
For subsets of cardinality at most three in the small-index routine,
the code tests the conclusion by the oracle; the paper proof of those
cases is Lemma~\ref{lem:three}.

\textbf{Local mixed-sumset audit.}
The program \code{audit\_local\_mixed.py} does not impose full
polynomiality. It generates $12{,}000$ inputs with $k\le2{,}500$,
finds the actual stabilizer modulus $g$ of the mixed sumset, and tests
only the local hypothesis $C_A(g)\ge C_B(g)$ used in
Remark~\ref{rem:mixedlocal}.
Exactly $9{,}840$ generated inputs satisfy this hypothesis and all pass.
The other $2{,}160$ are recorded as not meeting the hypothesis.
Of the passing inputs, $176$ have a nontrivial stabilizer and
$9{,}664$ have a trivial stabilizer.
An additional semigroup-oracle check is performed on $246$ of them.

\newpage
\textbf{Additional review audit.}
Version~2 includes a dedicated witness-construction check for
Lemma~\ref{lem:three}, on separated positive integer intervals, without
normalizing the overall greatest common divisor. It checks $124{,}927$
subsets on $3{,}651$ admissible interval pairs, including nonprimitive
triples and all remaining $v$-branches. Its scope and branch counts
are recorded in \path{review_v2/three_generator_audit.json} inside the
verification directory. This finite audit is supplementary; the interval
and representation arguments above establish the lemma.

\textbf{Reproducibility.}
The original supplementary archive and its scripts are preserved in
the source package. The following commands require Python~3.10 or
newer and its standard library:
\begin{quote}
\small
\begin{verbatim}
python verify_mixed_completion.py --k-max 16
  --random-cases 6000 --seed 202611103
  --out audit_main_rerun.json

python audit_local_mixed.py --cases 12000
  --seed 202611102 --out audit_local_rerun.json
\end{verbatim}
\end{quote}
Each displayed command is entered as one line.
Both commands were rerun while preparing this version.
Their output agrees with the supplied audit files in every field
except elapsed time.
All nine hashes in the supplied \code{SHA256.json} were verified;
the standalone proof draft and the draft inside the original archive
are byte-for-byte identical.
The supplement also records a counterexample to an unrestricted
double-sumset shortcut; that rejected assertion is not used in this paper.

\begin{thebibliography}{9}
\bibitem{GK}
Mona Gatzweiler and Christian Krattenthaler.
\emph{A positivity conjecture for a quotient of $q$-binomial coefficients}.
The Ramanujan Journal \textbf{69} (2026), article~13, 13~pp.
DOI: \href{https://doi.org/10.1007/s11139-025-01285-2}
{10.1007/s11139-025-01285-2}.
Preprint: \url{https://arxiv.org/abs/2502.06032}.

\bibitem{GLY}
Mona Gatzweiler, Fabi\'an Levic\'an-Santib\'a\~nez, and Atsuro Yoshida.
\emph{Cyclotomic generating functions, empty weighted complete
intersections and positivity}.
Preprint, arXiv:2603.22226v1 (2026).
\url{https://arxiv.org/abs/2603.22226v1}.

\bibitem{DeVos}
Matt DeVos.
\emph{A short proof of Kneser's addition theorem for abelian groups}.
Preprint, arXiv:1303.3539v1 (2013).
\url{https://arxiv.org/abs/1303.3539v1}.
\end{thebibliography}

\end{document}
